% Paragraph "Attended-name mover classification" of the Stage-3 plan (Section 1.4,
% after "Copying and Token-Preference Diagnostic"). This is the corrected version of the
% text pasted into Overleaf on 21 September 2026; the local
% Stage3_Head_Level_Characterization_revised.tex does not yet contain the paragraph.
% Change from the Overleaf version: the cross-check sentence now states the expected
% sign under the Stage-2 convention (the corrupted donor of a suppressor lowers the
% CORRUPTED-answer logit, i.e. negative Scope-F importance), instead of "lower the
% clean-answer logit". Implemented in pilot_v2/stage3_analyze.py (--attended-only).

\paragraph{Attended-name mover classification.}
The weight map describes the current token; the complementary question is what the head
does to the token it attends to. Using the saved test-block records of the common
families (attention argmax and the contextual output projection $o^h_j W_U$ over the
visible name tokens at the same position), select for each head the events whose argmax
falls on a name token and compute the contextual logit of that attended name's first
token relative to the mean over the other visible distinct names (both roles; last-token
anchor and same-role controls are sensitivity variants). A head is labeled a name mover
if this contrast is positive in at least 80\% of its events, and a negative name mover if
it is negative in at least 80\%, subject to the fixed support rule ($\geq 20$ events,
$\geq 10$ families); heads below support are reported as insufficient, not as neither.
The label is cross-checked against the sign of the saved 178-pair Scope-F importance:
under the Stage-2 convention (clean recipient, corrupted donor, sign fixed to the clean
answer) a head that suppresses the name it attends to yields a \emph{negative} Scope-F
importance, carried by a fall of the corrupted-answer logit, and a name mover a positive
one; a mismatch between the label and the sign of $I^F_h$ is reported as a finding rather
than resolved by relabeling, and $|I^F_h|<0.05$ is reported as ``no causal contrast to
check''. This classification uses only records already produced by the core capture and
requires no additional measurement. The candidates L23H15, L19H16 and L26H23 are checked
first but the classification is applied to the full inventory, including the random
controls.
